% SPDX-License-Identifier: Apache-2.0
% covers: Mmu, Mmu8
\datasheet{Mmu}{Sv32 address translation: a fetch TLB, a data TLB and a walker that never writes, answering in a register a cycle after each request.}

\dsfacts{\code{lib/parts/src/mmu.rs}}{\code{txhdl\_parts::mmu::Mmu}}%
{\code{//docs:vreteno}}

\subsection*{Function}

\code{Mmu<E>} translates a core's virtual addresses for Sv32, the
RISC-V page-based virtual memory of thirty-two bits (issue 1014). It
has two translation lookaside buffers of \code{E} entries each, one
for fetches and one for data, and one walker that fills them from the
page tables.

The core asks on \code{ireq} for a fetch and on \code{dreq} for a load
or a store, with a virtual address it already holds in a register. The
answer is a register too: on \code{ires} or \code{dres}, a cycle after
the request, a \code{Res} that says the access is allowed at a
physical address, or is a page fault, or an access fault. A miss
answers none of the three. The walker then reads the page table, a
word at a time, by sending the entry's address on \code{ptw} and
taking the entry, or the failure of its read, on \code{pte}. It fills
an entry, and the answer comes a cycle after the fill. The core holds
its request steady until one of the three comes. An answer belongs to
the request of the cycle before it.

The other inputs are \code{satp}, the mode the access is checked as,
\code{prv}, the two \code{mstatus} bits \code{sum} and \code{mxr}, and
\code{flush}, which empties both buffers.

\subsection*{Parameters}

\begin{center}\footnotesize
\begin{tabular}{@{}lp{0.7\textwidth}@{}}
\toprule
Parameter & Meaning \\
\midrule
\code{E} & The entries in each buffer, from 1 to 16. The tree uses 8, as
\code{Mmu8}, which is what the tables lower. \\
\bottomrule
\end{tabular}
\end{center}

\dstables{Mmu}

\subsection*{Behaviour}

\begin{itemize}
\item Translation is on when \code{satp.MODE} is set and \code{prv} is
not machine mode. Otherwise the address passes through unchecked, with
the same latency of one cycle, so the core has one rule.
\item Pages are four kilobytes. A leaf at the first level is a four
megabyte megapage, which compares the upper ten bits of the page
number only. Each buffer is fully associative and replaced in turn.
\item The checks are made on a hit: read, write and execute against the
access, the user bit against the mode and \code{sum}, and \code{mxr}
for a load from an execute-only page.
\item The walker never writes. An entry whose accessed bit is clear, or
whose dirty bit is clear on a store, is a page fault, and the kernel
sets the bit (Svade).
\item A flush empties both buffers. Address space identifiers are not
kept, which the specification allows, so \code{satp.ASID} is ignored.
A request in the cycle of a flush is not answered; the core holds it,
and it is answered from the emptied buffers, by a walk. A flush during
a walk drops what the walk finds.
\item When both ports miss at once, the data port is walked first.
\item A walk that ends in a fault answers for that page while its port
goes on asking for it, which the core does until it takes the trap,
and is not kept after.
\item Physical addresses are thirty-two bits. An entry or a
\code{satp} whose page number reaches past them is an access fault, as
is a table read that nothing answers.
\end{itemize}

\subsection*{Verification}

\begin{itemize}
\item \code{translate} in \code{mmu.rs} is the reference: the
specification's walk, with no buffer. In \code{mmu::tests}, in
\code{//lib/parts:parts\_test}, the unit is held to it against tables
in a memory that answers the walker after a few cycles:
\code{bare\_and\_machine\_mode}\allowbreak\code{\_pass\_the\_address}\allowbreak\code{\_through\_a\_cycle\_later};
\code{a\_page\_is\_walked\_once}\allowbreak\code{\_and\_then\_hits\_a\_cycle\_later};
\code{a\_megapage\_maps}\allowbreak\code{\_four\_megabytes\_from\_one\_entry};
\code{every\_fault}\allowbreak\code{\_is\_the\_specifications};
\code{an\_address\_past}\allowbreak\code{\_thirty\_two\_bits}\allowbreak\code{\_or\_nowhere}\allowbreak\code{\_is\_an\_access\_fault};
\code{a\_fault\_is\_not\_kept}\allowbreak\code{\_after\_its\_request};
\code{a\_flush\_empties\_both\_tlbs};
\code{a\_flush\_during\_a\_walk}\allowbreak\code{\_drops\_what\_it\_finds};
\code{a\_request\_in\_the\_cycle}\allowbreak\code{\_of\_a\_flush\_is\_answered\_after\_it};
\code{entries\_are\_replaced\_in\_turn};
\code{the\_data\_port\_walks\_first}; and
\code{random\_requests}\allowbreak\code{\_agree\_with\_the\_reference}, six hundred random
requests over random tables, on both ports and in every mode.
\code{the\_eight\_entry\_unit\_lowers} writes \code{Mmu8}'s Verilog.
\item \code{ex\_mmu} runs \code{Mmu8} the way a kernel meets it, and
the build simulates its netlist against that run under nvc and
Verilator: \code{//docs:mmu\_sim\_mmu8\_tb\_test} and
\code{//docs:mmu\_vsim\_test}.
\item In the core, \code{//cpu/vreteno:lockstep\_test} runs Sv32
directed and over random page tables against the reference model
(\code{sv32\_translates\_and\_faults}\allowbreak\code{\_where\_the\_tables\_say},
\code{random\_programs\_under\_paging}), and Linux runs under it on the
board (issue 279).
\item \code{//lib/parts:mmu\_pnr}, which is manual, places and routes
\code{Mmu8} alone, out of context, against the core's ten nanosecond
clock: it meets it with 2.440\,ns to spare, in 493 LUTs and 919
flip-flops (\code{//docs:vreteno}).
\end{itemize}

\subsection*{Used by}

\code{Mmu8} in \code{Hart}, the core and its unit, in
\code{cpu/vreteno/src/hart.rs}, and so in every design that holds a
hart: the board, the flagship and \code{Pair}. \code{Mmu8} in the
example \code{ex\_mmu}.

\subsection*{Limits}

No address space identifiers. No hardware update of the accessed and
dirty bits. The buffers are held in flip-flops, not block RAM, so
\code{E} stays small. Physical addresses are thirty-two bits, with no
Sv32 thirty-four bit physical space.
